% Dedicatorias por autor (2026-10-04), redactadas a partir de las respuestas de cada autor.
\chapter*{Dedicatoria}

\begingroup
\setstretch{1.15}
\setlength{\parindent}{0pt}
\setlength{\parskip}{0.8em}

\begin{flushright}
\begin{minipage}{0.78\linewidth}
\raggedleft\itshape
«Si he logrado ver más lejos, ha sido porque he subido a hombros de gigantes».\par
\upshape ---Isaac Newton, carta a Robert Hooke (1676)
\end{minipage}
\end{flushright}

\vspace{1.2em}

A Dios, primero que todo, por cada bendición y cada oportunidad que puso en mi camino. A mis padres, Georgina y Reinaldo, por su sacrificio, por su apoyo económico y moral y por el amor incondicional que siempre me han brindado. A mis hermanos, Cristian, Duben y Camilo, por su lealtad y por haber estado siempre presentes. A mi nona Toña y a mi abuelita Edelmira, que en paz descanse. A mis amigos Daniel Pinto, Gabriel Castillo, Juan Camilo Gómez, Camilo Gamboa, Brayan Quintero, Jeferson Acevedo, Luis Mora, César Camargo, Andrés León y Juan Camilo Cuadrado, y en especial a Gabriel y a Juan Camilo Gómez, por su apoyo incondicional, sus consejos y una amistad inquebrantable. Y a don Gonzalo y doña Dolly, mis vecinos.

\begin{flushright}\textit{\AutorUno}\end{flushright}

\vspace{2em}

A Dios. A mis padres, Luis Pinto y Saidy Ortega, y a mis hermanos, Juan e Isabella Pinto, por sus sacrificios y sus consejos. A mis abuelos en el cielo, Ciro Antonio Pinto Torres y Elizabeth Rizzo, y a mi abuelita María Hernández. A mis amigos Reinaldo, mi compañero de tesis, Sharon Ortega, Henry Salas, Otto Chacón, Marlon Cely, Sergio Alvarado y Juan Urrego. Porque me enseñaron que ninguna parte funciona aislada: este logro no es el esfuerzo de un elemento individual, sino el resultado de un sistema perfecto de apoyo, amor y guía que fluyó siempre a mi favor.

\begin{flushright}\textit{\AutorDos}\end{flushright}

\endgroup
